\documentclass[12pt]{article}
\usepackage{graphicx} % Required for inserting images
\usepackage{mathtools}
\usepackage{graphicx}
\usepackage{subcaption}
\usepackage{float}
\usepackage{booktabs}
\usepackage{hyperref}
\usepackage[
    top=1in,
    bottom=1in,
    left=1in,
    right=1in
]{geometry}
\usepackage{amsmath}
\usepackage{amssymb}
\title{\Huge \textbf{photonForge}}

\begin{document}

\maketitle

\section{Abstract}
photonForge is our attempt to make the phenomenon of black holes — long familiar to physicists but difficult for students to intuitively visualise — computationally and visually accessible. In this project, we develop a computational framework for rendering physically accurate black hole images using relativistic ray tracing based on the Schwarzschild and Kerr spacetime metrics, and extend this beyond conventional visualisation by incorporating machine learning techniques into the computational pipeline.

\section{Participants Info}

\begin{itemize}
    \item \texttt{Anshika Sharma}\\
    University of Allahabad\\
    \href{mail :usernameyjs.1972@gmail.com}{usernameyjs.1972@gmail.com}

    \item \texttt{Manya Mishra}\\
    University of Allahabad\\
    \href{mail :manyamishrasmc6002@gmail.com}{manyamishrasmc6002@gmail.com}

    \item \texttt{Vagisha Yadav}\\
    University of Allahabad\\
    \href{mail :vagisha.allduniv@gmail.com}{vagisha.allduniv@gmail.com}
\end{itemize}




\section{Competition Category}
Category 1- Undergraduate (UG)– Foundations of Computational Physics (Orbital Dynamics)
\newpage
\section{Statement of Problem}
Real time rendering of physically accurate gravitational lensing around Schwarzchild and kerr black holes is computationally intractable using per-pixel numerical geodesic integration at interactive frame rates. This project formulates black-hole images rendering as a learned function approximation problem by replacing per pixel RK4 geodesic integration with a neural network surrogate suitable for GPU shader inference, trained on a physically accurate simulation pipeline.

\section{Introduction}
The universe remains one of the most enduring subjects of inquiry in physics, engaging thinkers from Albert Einstein to Stephen Hawking across the past century. Among its many mysteries, black holes remain especially elusive: by definition, they absorb all light and radiation that crosses their event horizon, and so can never be observed directly — only inferred through their gravitational effects on their surroundings. Their existence emerges as a solution of Einstein's general theory of relativity (1915), first derived explicitly by Karl Schwarzschild in 1916.
The first mathematically grounded image of a black hole was produced by physicist Jean-Pierre Luminet in 1979, computed on an early mainframe and hand-rendered using India ink — decades before modern scientific computing tools existed. Luminet noted that his results should extend to a range of black hole systems, including the supermassive black hole suspected at the core of the elliptical galaxy M87. In April 2019, the Event Horizon Telescope Collaboration released the first direct telescopic image of the shadow of M87* and its surrounding accretion disk — a striking confirmation of predictions made forty years earlier.

A black hole is a region of spacetime whose gravitational field is strong enough that no signal, including light, can escape from within its event horizon. Its observable appearance to a distant observer is governed entirely by the paths that photons trace through the curved spacetime surrounding it — null geodesics — rather than by any emitted light from the black hole itself. This project computationally reconstructs that appearance by directly solving the equations governing these photon paths, for both a non-rotating (Schwarzschild) and a rotating (Kerr) black hole. It renders the resulting image including the effects of an orbiting accretion disk. 

\newpage
\section{\textbf{Literature Review and Formulation of Problem}}
The rendering of astrophysical images from the vicinity of black holes has long depended on numerically tracing the paths of photons through curved spacetime. Beyond its scientific applications, physically accurate rendering of gravitational lensing has proven valuable for public science communication and education — most notably in the Double Negative Gravitational Renderer (DNGR)\cite{james2015}, developed to depict a spinning black hole for the film Interstellar (James et al. 2015), which demonstrated that IMAX-quality, flicker-free imagery of a Kerr black hole could be produced using ray-bundle propagation techniques rather than single-ray integration. As interactive and real-time applications of such visualizations become more common, in education, exploration tools, and browser-based astronomical visualization, the computational cost of generating each frame becomes a central bottleneck, motivating the search for faster alternatives to direct numerical ray-tracing.

The general relativistic ray-tracing problem has been addressed by numerous codes to date (e.g., Vincent et al. 2011\cite{vincent2011}; Chan et al. 2013\cite{chan2013}; James et al. 2015\cite{james2015}), following two broad computational strategies. In one approach, applicable only to integrable spacetimes such as Kerr, conserved quantities (energy, angular momentum, and the Carter constant) are used to reduce the geodesic equations to a lower, semi-analytic order, permitting large integration steps at the cost of generality (e.g., Dexter and Agol 2009\cite{dexter2009}). In the other, more general approach, the second-order geodesic equations are integrated directly via numerical schemes, most commonly a 4th-order Runge-Kutta (RK4) integrator, allowing extension to non-Kerr and numerically computed metrics but at substantially higher per-photon computational cost. GYOTO (Vincent et al. 2011\cite{vincent2011}) exemplifies this direct-integration approach on the CPU, using a Hamiltonian RK4 formulation with explicit correction of the conserved quantities at each step, and supports both analytic Kerr metrics and arbitrary numerical (3+1) spacetimes, though at a reported cost of minutes per rendered frame. GRay (Chan et al. 2013\cite{chan2013}) applies the same class of RK4 integration but restructures it for massively parallel GPU execution via NVIDIA's CUDA framework, achieving up to three orders of magnitude speedup over equivalent serial CPU codes and enabling large-scale parameter studies of photon-ring and shadow properties across black hole spin and inclination. DNGR (James et al. 2015\cite{james2015}) instead propagates elliptical ray bundles rather than individual rays, trading a more involved formulation for smoother, flicker-free imagery under a moving camera. Despite their differences in platform and formulation, all three approaches share a common limitation: every rendered pixel requires the underlying differential equations to be integrated from scratch, with no mechanism for amortizing this cost across pixels, frames, or black hole parameters.

More recently, neural network surrogates have been proposed to replace direct geodesic integration entirely, following a broader trend in scientific machine learning toward learned operators that approximate expensive simulation pipelines with a single forward pass. GravLensX (Sun et al. 2025\cite{sun2025}) trains a neural network to fit the null-geodesic mapping around single and superposed Kerr black holes, using the trained model in place of direct integration to render optically thin accretion disks, and reports 15–26× speedup over conventional ray-tracing depending on scene complexity. A parallel line of work embeds the governing differential equations directly into the training objective rather than learning purely from precomputed data: physics-informed neural networks (PINNs; Raissi et al. 2019\cite{raissi2019}) have been shown capable of reconstructing relativistic orbital dynamics directly from observational data, including the precessing orbit of the S2 star around Sgr A* (Eyama and Masada 2025\cite{eyama2025}), though such approaches are typically trained per-instance and are known to struggle with the stiff, branch-sensitive dynamics characteristic of photon trajectories near the photon sphere. Hybrid physics-informed operator architectures (e.g., PINO-MBD; Ding et al. 2021) that combine parameter-space generalization with physics-constrained training remain largely unexplored for geodesic ray-tracing specifically.

Taken together, existing work establishes that direct numerical integration, whether CPU-based and physically comprehensive (GYOTO), GPU-parallel and fast (GRay), or reformulated for smooth interactive imagery (DNGR), remains fundamentally limited by its per-pixel computational cost, while the nascent neural-surrogate literature (GravLensX) demonstrates that this cost can be substantially reduced but has not yet been targeted at in-shader, real-time deployment. This project addresses that gap by training a multi-task neural network surrogate — a shared MLP trunk with dedicated outcome-classification, disk-intersection, and background-direction regression heads — on RK4-integrated ground-truth geodesic data, for deployment as either in-shader GLSL inference or a baked lookup texture within a real-time WebGL2 rendering pipeline.

The visualization of black holes and the phenomenon of gravitational lensing are often inaccessible to the general public due to the complexity of general relativity and the high computational demands of existing research-grade simulators. Because widely available simplified visualizations frequently misrepresent the underlying physics in the interest of accessibility, there is a critical need for a computational framework that can accurately compute and visualize these phenomena without sacrificing educational clarity. To address this gap, photonForge outlines a framework that numerically studies the dynamics of photon trajectories in curved spacetime, extending the concepts of classical orbital dynamics, such as reliance on conserved quantities and effective-potential-based trajectory analysis to massless particles. 

Physically and mathematically, this requires computing the trajectories of freely propagating photons along null geodesics ($ds^{2}=0$) within both the non-rotating Schwarzschild and rotating Kerr spacetime metrics. Initially, these trajectories are computed by numerically integrating the equations using a fourth-order Runge-Kutta (RK4) integrator. However, because directly solving these coupled differential equations per pixel creates a severe performance bottleneck, the computational challenge is ultimately reformulated as a machine learning optimization task. To achieve real-time rendering speeds, a multi-task neural network surrogate is developed to map initial photon impact parameters directly to final physical observables—such as minimum radial distance, deflection angle, and event horizon capture state—effectively bypassing the heavy computational cost of traditional ray tracing while preserving strict relativistic accuracy.
\newpage
\section{\textbf{THEORY}}
\subsection{Einstein Field Equations}

Einstein's Field Equations relate the curvature of spacetime to the distribution of matter and energy within it. They form the foundation of General Relativity and may be written as

\begin{equation}
R_{\mu\nu} - \frac{1}{2}g_{\mu\nu}R
=
\frac{8\pi G}{c^4}T_{\mu\nu},
\end{equation}

or equivalently,

\begin{equation}
G_{\mu\nu}
=
\frac{8\pi G}{c^4}T_{\mu\nu},
\end{equation}

where the Einstein tensor is defined as

\begin{equation}
G_{\mu\nu}
=
R_{\mu\nu}
-
\frac{1}{2}g_{\mu\nu}R.
\end{equation}

The quantities appearing in the field equations are:

\begin{itemize}
    \item $G_{\mu\nu}$: Einstein tensor, representing the curvature of spacetime.
    \item $g_{\mu\nu}$: Metric tensor, which defines the geometry of spacetime.
    \item $R_{\mu\nu}$: Ricci curvature tensor.
    \item $R$: Ricci scalar, obtained by contracting the Ricci tensor,
    \[
    R = g^{\mu\nu}R_{\mu\nu}.
    \]
    \item $T_{\mu\nu}$: Stress-energy tensor, describing the distribution of mass, energy, momentum, and pressure.
    \item $G$: Newton's gravitational constant.
    \item $c$: Speed of light in vacuum.
\end{itemize}

The Ricci tensor is obtained by contracting the Riemann curvature tensor,

\begin{equation}
R_{\mu\nu}
=
R^{\lambda}_{\ \mu\lambda\nu},
\end{equation}

where the Riemann curvature tensor is given by

\begin{equation}
R^{\rho}_{\ \sigma\mu\nu}
=
\partial_{\mu}\Gamma^{\rho}_{\nu\sigma}
-
\partial_{\nu}\Gamma^{\rho}_{\mu\sigma}
+
\Gamma^{\rho}_{\mu\lambda}\Gamma^{\lambda}_{\nu\sigma}
-
\Gamma^{\rho}_{\nu\lambda}\Gamma^{\lambda}_{\mu\sigma}.
\end{equation}

The Christoffel symbols are defined in terms of the metric as

\begin{equation}
\Gamma^{\alpha}_{\beta\gamma}
=
\frac{1}{2}
g^{\alpha\sigma}
\left(
\frac{\partial g_{\sigma\beta}}{\partial x^{\gamma}}
+
\frac{\partial g_{\sigma\gamma}}{\partial x^{\beta}}
-
\frac{\partial g_{\beta\gamma}}{\partial x^{\sigma}}
\right).
\end{equation}

The Einstein Field Equations constitute a system of ten coupled, nonlinear, second-order partial differential equations. Solving these equations for a given stress-energy distribution determines the metric tensor $g_{\mu\nu}$ and therefore the geometry of spacetime.

\subsection{Schwarzschild Metric}

The Einstein Field Equations are highly non-linear and, in general, difficult to solve analytically. However, for the special case of a static, spherically symmetric, electrically neutral, and non-rotating mass in vacuum, Karl Schwarzschild obtained the first exact solution to the field equations in 1916. This solution, known as the Schwarzschild metric, describes the geometry of spacetime surrounding such a mass. It replaces the Euclidean description of distance with a spacetime interval that incorporates the effects of gravity.
The corresponding metric tensor is

\begin{equation}
g_{\mu\nu}
=
\begin{pmatrix}
-\left(1-\dfrac{2M}{r}\right) & 0 & 0 & 0\\[0.8em]
0 & \left(1-\dfrac{2M}{r}\right)^{-1} & 0 & 0\\[0.8em]
0 & 0 & r^2 & 0\\[0.8em]
0 & 0 & 0 & r^2\sin^2\theta
\end{pmatrix},
\end{equation}

where the coordinates are ordered as

\[
x^\mu = (t,\, r,\, \theta,\, \phi).
\]

Since all off-diagonal components vanish, the Schwarzschild metric is diagonal. The diagonal entries describe how spacetime measures intervals in the temporal, radial, and angular directions.

The line element is hence given by

\begin{equation}
ds^2
=
-\left(1-\frac{2M}{r}\right)dt^2
+
\left(1-\frac{2M}{r}\right)^{-1}dr^2
+
r^2d\theta^2
+
r^2\sin^2\theta\,d\phi^2,
\end{equation}

where $M$ is the mass of the gravitating object, $r$ is the radial coordinate, $t$ is the time coordinate, and $(\theta,\phi)$ are the angular coordinates.

Each term in the metric has a distinct physical interpretation:

\begin{itemize}
    \item $\displaystyle -\left(1-\frac{2M}{r}\right)dt^2$: Represents \textbf{gravitational time dilation}; how the passage of time slows down as one approaches the black hole

    \item $\displaystyle \left(1-\frac{2M}{r}\right)^{-1}dr^2$: Represents the \textbf{radial component} of the spacetime interval; accounts for the stretching of radial distances due to spacetime curvature, causing radial lengths to appear larger than they would in flat spacetime.

    \item $\displaystyle r^2d\theta^2 + r^2\sin^2\theta\,d\phi^2$: Represents the \textbf{angular geometry} of spacetime. Owing to spherical symmetry, every surface of constant radius remains a perfect sphere with surface area $4\pi r^2$, even though the radial geometry is curved.

\subsection{Why Photons Follow Null Geodesics}

In general relativity, the spacetime interval between two infinitesimally separated events is given by

\[
ds^2 = g_{\mu\nu}\,dx^\mu dx^\nu.
\]
\begin{center}
where $ds^2$ is the spacetime interval,\\
$g_{\mu\nu}$ is the metric tensor,\\
$dx^\mu, dx^\nu$ are infinitesimal changes in spacetime coordinates,\\
and $\mu,\nu$ are indices running over the four coordinates $(0,1,2,3)$.
\end{center}

The nature of a worldline is determined by the value of \(ds^2\). A timelike worldline satisfies \(ds^2 < 0\), a null worldline satisfies \(ds^2 = 0\), and a spacelike worldline satisfies \(ds^2 > 0\), for the metric signature \((-+++)\).

Consider first flat spacetime, whose metric is

\[
ds^2 = -c^2dt^2 + dx^2 + dy^2 + dz^2.
\]

A photon travels at the speed of light, \(c\). Hence, during an infinitesimal time interval \(dt\), the spatial distance travelled by the photon is

\[
dl = c\,dt.
\]

Therefore,

\[
dx^2 + dy^2 + dz^2 = dl^2 = c^2dt^2.
\]

Substituting this into the spacetime interval gives

\[
ds^2 = -c^2dt^2 + c^2dt^2 = 0.
\]

Thus, the worldline of a photon is a \emph{null curve}:

\[
\boxed{ds^2=0}.
\]

This condition does not imply that the photon is stationary or that its spatial and temporal components vanish. Rather, the temporal and spatial contributions to the spacetime interval exactly cancel because the photon propagates at \(c\).

In curved spacetime, the interval is instead determined by the spacetime metric,

\[
ds^2 = g_{\mu\nu}\,dx^\mu dx^\nu.
\]

Since photons remain massless and propagate locally at the speed of light, their worldlines continue to satisfy

\[
\boxed{ds^2=0}.
\]
However, the condition \(ds^2=0\) alone only establishes that the path is null; it does not by itself determine the trajectory. A freely propagating photon follows a geodesic, meaning that its covariant acceleration vanishes:

\[
\frac{D^2x^\mu}{d\lambda^2}=0,
\]

where \(\lambda\) is an affine parameter along the photon trajectory. Combining the null condition with the geodesic equation therefore gives the equations governing photon propagation in curved spacetime.

Consequently, photons follow \emph{null geodesics}: they follow geodesics because they are freely propagating particles in curved spacetime, and those geodesics are null because photons have zero rest mass and travel at the speed of light.
    \subsection{Derivation of the Null Geodesic Equations}

The trajectory of a photon is described by null geodesics of the Schwarzschild spacetime. Beginning with the Schwarzschild line element,

\begin{equation}
ds^2=
-\left(1-\frac{2M}{r}\right)dt^2
+\left(1-\frac{2M}{r}\right)^{-1}dr^2
+r^2d\theta^2
+r^2\sin^2\theta\,d\phi^2,
\end{equation}

we restrict the motion to the equatorial plane,

\begin{equation}
\theta=\frac{\pi}{2},
\qquad
d\theta=0,
\end{equation}

so that the metric becomes

\begin{equation}
ds^2=
-\left(1-\frac{2M}{r}\right)dt^2
+\left(1-\frac{2M}{r}\right)^{-1}dr^2
+r^2d\phi^2.
\end{equation}

Introducing an affine parameter $\lambda$, the geodesic Lagrangian is

\begin{equation}
\mathcal{L}
=
\frac12
g_{\mu\nu}
\frac{dx^\mu}{d\lambda}
\frac{dx^\nu}{d\lambda},
\end{equation}

substituting from the metric,

\begin{equation}
\mathcal{L}
=
\frac12
\left[
-\left(1-\frac{2M}{r}\right)
\left(\frac{dt}{d\lambda}\right)^2
+
\left(1-\frac{2M}{r}\right)^{-1}
\left(\frac{dr}{d\lambda}\right)^2
+
r^2
\left(\frac{d\phi}{d\lambda}\right)^2
\right].
\end{equation}

Since the coordinates $t$ and $\phi$ do not appear explicitly in the Lagrangian, they are cyclic coordinates. Applying the Euler--Lagrange equations gives two conserved quantities,

\begin{equation}
\left(1-\frac{2M}{r}\right)
\frac{dt}{d\lambda}
=
E,
\end{equation}

and

\begin{equation}
r^2
\frac{d\phi}{d\lambda}
=
L,
\end{equation}

where $E$ and $L$ are constants corresponding to the conserved energy and angular momentum of the photon.

Hence,

\begin{equation}
\frac{dt}{d\lambda}
=
\frac{E}{1-\frac{2M}{r}},
\end{equation}

and

\begin{equation}
\frac{d\phi}{d\lambda}
=
\frac{L}{r^2}.
\end{equation}

Since photons travel on null geodesics,

\begin{equation}
ds^2=0,
\end{equation}

which implies

\begin{equation}
-\left(1-\frac{2M}{r}\right)
\left(\frac{dt}{d\lambda}\right)^2
+
\left(1-\frac{2M}{r}\right)^{-1}
\left(\frac{dr}{d\lambda}\right)^2
+
r^2
\left(\frac{d\phi}{d\lambda}\right)^2
=
0.
\end{equation}

Substituting the conserved quantities yields

\begin{equation}
-
\left(1-\frac{2M}{r}\right)
\left(\frac{E}{1-\frac{2M}{r}}\right)^2
+
\left(1-\frac{2M}{r}\right)^{-1}
\left(\frac{dr}{d\lambda}\right)^2
+
r^2
\left(\frac{L}{r^2}\right)^2
=
0.
\end{equation}

Multiplying throughout by $\left(1-\frac{2M}{r}\right)$ and simplifying gives

\begin{equation}
\boxed{
\left(\frac{dr}{d\lambda}\right)^2
=
E^2
-
\left(1-\frac{2M}{r}\right)
\frac{L^2}{r^2}
}
\end{equation}

Therefore, the null geodesic equations governing photon motion in Schwarzschild spacetime are

\begin{align}
\frac{dt}{d\lambda}
&=
\frac{E}{1-\frac{2M}{r}},\\[0.5em]
\frac{dr}{d\lambda}
&=
\pm
\sqrt{
E^2-
\left(1-\frac{2M}{r}\right)
\frac{L^2}{r^2}
},\\[0.5em]
\frac{d\phi}{d\lambda}
&=
\frac{L}{r^2}.
\end{align}
\subsection{Null Condition}

\textit{(Path of the photon)}

\begin{equation}
ds^2 = 0
\end{equation}

\begin{equation}
-\left[1-\frac{2M}{r}\right]
\left(\frac{dt}{d\lambda}\right)^2
+
\left[1-\frac{2M}{r}\right]^{-1}
\left(\frac{dr}{d\lambda}\right)^2
+
r^2
\left(\frac{d\phi}{d\lambda}\right)^2
=0
\end{equation}

Substituting,

\begin{equation}
\frac{dt}{d\lambda}
=
\frac{E}{1-\frac{2M}{r}}
\end{equation}

\begin{equation}
\frac{d\phi}{d\lambda}
=
\frac{L}{r^2}
\end{equation}

Therefore,

\begin{equation}
-E^2
+
\left(\frac{dr}{d\lambda}\right)^2
+
\left[1-\frac{2M}{r}\right]
\left[\frac{L^2}{r^2}\right]
=0
\end{equation}

Hence,

\begin{equation}
\left(\frac{dr}{d\lambda}\right)^2
=
E^2
-
\left[1-\frac{2M}{r}\right]
\left[\frac{L^2}{r^2}\right]
\end{equation}

Where the effective potential is

\begin{equation}
V_{\mathrm{eff}}(r)
=
\left[\frac{L^2}{r^2}\right]
\left[1-\frac{2M}{r}\right]
\end{equation}

Therefore,

\begin{equation}
\left(\frac{dr}{d\lambda}\right)^2
=
E^2-V_{\mathrm{eff}}(r)
\end{equation}


\subsection{Photon Sphere}

For the photon sphere, $V_{\mathrm{eff}}(r)$ has a maximum. Therefore,

\begin{equation}
\frac{dV_{\mathrm{eff}}}{dr}=0
\end{equation}

The effective potential can be written as

\begin{equation}
V_{\mathrm{eff}}
=
L^2
\left[
\frac{1}{r^2}
-
\frac{2M}{r^3}
\right]
\end{equation}

Differentiating,

\begin{equation}
\frac{dV_{\mathrm{eff}}}{dr}
=
L^2
\left[
-\frac{2}{r^3}
+
\frac{6M}{r^4}
\right]
=0
\end{equation}

Multiplying by $r^4$,

\begin{equation}
-2r+6M=0
\end{equation}

Therefore,

\begin{equation}
r=3M
\end{equation}

Thus, the radius of the photon sphere is

\begin{equation}
\boxed{r_{\mathrm{ph}}=3M}
\end{equation}

At $r=3M$, the effective potential is

\begin{equation}
V_{\mathrm{eff}}
=
L^2
\left[
\frac{1}{9M^2}
-
\frac{2M}{27M^3}
\right]
\end{equation}

\begin{equation}
V_{\mathrm{eff}}
=
\frac{L^2}{M^2}
\left[
\frac{1}{9}
-
\frac{2}{27}
\right]
\end{equation}
\newpage
\subsection{Kerr Metric}

The Schwarzschild solution describes a static, spherically symmetric, non-rotating mass. However, astrophysical black holes are generally expected to possess angular momentum. The corresponding exact solution to the Einstein Field Equations for a stationary, axisymmetric, electrically neutral, rotating mass was obtained by Roy Kerr in 1963. This solution, known as the Kerr metric, describes the geometry of spacetime surrounding a rotating black hole.

In Boyer--Lindquist coordinates, the corresponding metric tensor is
The corresponding metric tensor is

\begin{equation}
g_{\mu\nu}
=
\begin{pmatrix}
-\left(1-\dfrac{2Mr}{\Sigma}\right)
&
0
&
0
&
-\dfrac{2Mar\sin^2\theta}{\Sigma}
\\[0.8em]
0
&
\dfrac{\Sigma}{\Delta}
&
0
&
0
\\[0.8em]
0
&
0
&
\Sigma
&
0
\\[0.8em]
-\dfrac{2Mar\sin^2\theta}{\Sigma}
&
0
&
0
&
\left(
r^2+a^2+\dfrac{2Ma^2r\sin^2\theta}{\Sigma}
\right)\sin^2\theta
\end{pmatrix}.
\end{equation}

where the coordinates are ordered as

\begin{equation}
x^\mu=(t,r,\theta,\phi).
\end{equation}

The quantities $\Sigma$ and $\Delta$ are defined as

\begin{equation}
\Sigma=r^2+a^2\cos^2\theta,
\end{equation}

and

\begin{equation}
\Delta=r^2-2Mr+a^2.
\end{equation}

Here, $M$ represents the mass of the black hole and $a$ represents its specific angular momentum, defined by

\begin{equation}
a=\frac{J}{M},
\end{equation}

where $J$ is the angular momentum of the black hole.

Unlike the Schwarzschild metric, the Kerr metric contains non-zero off-diagonal components, specifically the $t\phi$ and $\phi t$ components. These arise due to the rotation of the black hole and represent the coupling between the temporal and azimuthal coordinates.

The line element is hence given by
\begin{equation}
\begin{aligned}
ds^2 ={}&
-\left(1-\frac{2Mr}{\Sigma}\right)dt^2
-\frac{4Mar\sin^2\theta}{\Sigma}\,dt\,d\phi
\\
&+\frac{\Sigma}{\Delta}\,dr^2
+\Sigma\,d\theta^2
\\
&+\left(
r^2+a^2+
\frac{2Ma^2r\sin^2\theta}{\Sigma}
\right)
\sin^2\theta\,d\phi^2.
\end{aligned}
\end{equation}

Each term in the metric has a distinct physical interpretation:


\begin{itemize}

\item[$\bullet$] $\displaystyle -\left(1-\frac{2Mr}{\Sigma}\right)dt^2$: Represents \textbf{the temporal component} of the spacetime interval; describes the gravitational effects on the passage of time near the rotating black hole.

\item[$\bullet$] $\displaystyle -\frac{4Mar\sin^2\theta}{\Sigma}\,dt\,d\phi$: Represents \textbf{the frame-dragging component} of the spacetime interval; arises from the rotation of the black hole and describes the coupling between time and azimuthal motion.

\item[$\bullet$] $\displaystyle \frac{\Sigma}{\Delta}\,dr^2$: Represents \textbf{the radial component} of the spacetime interval; accounts for the modification of radial distances due to the curvature of spacetime.

\item[$\bullet$] $\displaystyle \Sigma\,d\theta^2$: Represents \textbf{the polar component} of the spacetime interval; describes the geometry associated with motion in the polar direction.

\item[$\bullet$] $\displaystyle\left(r^2+a^2+\frac{2Ma^2r\sin^2\theta}{\Sigma}\right)\sin^2\theta\,d\phi^2$:
: Represents \textbf{the azimuthal component} of the spacetime interval; describes distances associated with motion around the axis of rotation.

\end{itemize}

In the limit $a\rightarrow0$, the Kerr metric reduces to the Schwarzschild metric. Thus, the Schwarzschild solution can be regarded as the non-rotating limit of the more general Kerr solution.

\section{\textbf{Methodology, Implementation and \\testing of the code}}
The Schwarzschild and Kerr solutions of the Einstein Field Equations provide the spacetime metrics from which the null geodesic equations governing photon propagation are formulated. These equations are expressed as a system of coupled first-order ordinary differential equations parameterized by the affine parameter $\lambda$, and are solved independently for each pixel of the rendering screen using a fourth-order Runge–Kutta (RK4) integrator. The resulting photon trajectories are evaluated against a set of predefined termination conditions corresponding to the event horizon, the accretion disk, and escape to the background celestial sphere. Photons that cross the event horizon are considered captured and rendered as black; photons that intersect the accretion disk are assigned colour values derived from its radiation profile; and photons that escape the gravitational field are mapped onto a background star field following their net angular deflection. The aggregation of these pixel-wise trajectory outcomes constitutes the final rendered image of the black hole.

To implement this general relativistic ray-tracing and orbital simulation framework, a modular, object-oriented software architecture was developed in Python, complemented by a high-performance WebGL2 GPU rendering engine for real-time visualization. The core computational pipeline is structured as follows:

\subsection{Implementation Framework}

\subsubsection{\textbf{Metric Abstraction and Coordinate Systems:}}
Spacetime geometry is represented through an abstract base class \texttt{Metric}, which defines a common interface for coordinate transformations and evaluation of the 
metric tensor $g_{\mu\nu}$. This base class is specialized into \small\texttt{SchwarzschildMetric}, representing static, spherically symmetric spacetimes, and 
\small\texttt{KerrMetric}, representing rotating, axially symmetric spacetimes expressed in Boyer--Lindquist coordinates.
\medskip

\subsubsection{\textbf{Numerical Christoffel Symbol Solver:}}
Rather than hard-coding coordinate-specific equations of motion, the framework computes the Christoffel symbols $\Gamma^{\mu}_{\alpha\beta}$ dynamically using finite difference approximations of the metric tensor derivatives:
\begin{equation}
\Gamma^{\mu}_{\alpha\beta}
=
\frac{1}{2}g^{\mu\gamma}
\left(
\frac{\partial g_{\gamma\beta}}{\partial x^\alpha}
+
\frac{\partial g_{\alpha\gamma}}{\partial x^\beta}
-
\frac{\partial g_{\alpha\beta}}{\partial x^\gamma}
\right).
\end{equation}

This numerical formulation of the affine connection allows arbitrary metric definitions to be resolved dynamically within the geodesic solver, without requiring metric-specific derivations.

\medskip

\subsubsection{\textbf{Geodesic Solver and Adaptive RK4 Integration:}}
The second-order geodesic differential equations are reduced to a system of eight coupled first-order ordinary differential equations in the 4-position $x^\mu$ and 4 velocity $u^\mu$, and integrated using a fourth-order Runge--Kutta scheme (\texttt{RK(4) Integrator}).

To mitigate numerical instability and step-size accumulation in regions of high spacetime curvature, such as near the event horizon, a curvature-adaptive step-size controller (\texttt{Adaptive RK4}) dynamically scales the integration step $h$ as a function of the radial coordinate $r$:

\begin{equation}
h(r)
=
h_{\min}
+
\min\left[
\frac{r-r_{\mathrm{photon}}}{17.0},
1.0
\right]
\left(
h_{\max}-h_{\min}
\right).
\end{equation}

\medskip

\subsubsection{\textbf{Local Orthonormal Tetrads and Ray Generation:}}
To simulate physically realistic camera views, local observer reference frames are constructed via \texttt{ObserverTetrad} and \texttt{KerrObserverTetrad}. These convert the rectangular pixel grid of the camera's image plane into initial coordinate velocities $u^\mu$ oriented toward the black hole.

\medskip

\subsubsection{\textbf{Visual and Accretion Disk Physics:}}
The reference ray tracers (\small\texttt{schwarzschild\_ray\_tracer.py} and\\ \texttt{kerr\_ray\_tracer.py}) integrate photon trajectories backward in time from the  
observer's screen to determine intersection with the accretion disk or capture by the event horizon. The accretion disk is shaded according to a Keplerian velocity  
profile, incorporating relativistic Doppler beaming and gravitational redshift.
\medskip

\subsubsection{\textbf{Real-Time GPU WebGL2 Engine:}}
An interactive web application employs WebGL2 shaders to perform real-time ray-marching, enabling rapid, interactive exploration of black hole parameters---mass $M$ and spin $a$---at sustained frame rates of 60 FPS.


\section{Testing of Code}

\subsection{Additional Validation: Trajectory Classes and Relativistic Effects}

Beyond the conservation-law and orthonormality checks described above, the geodesic solver was further validated against a series of analytically and qualitatively known relativistic phenomena, spanning both the null (photon) and timelike (massive particle) sectors of the Schwarzschild and Kerr metrics. Each validation targets a specific, independently verifiable physical prediction, providing cumulative confidence that the numerical framework correctly reproduces general relativistic dynamics across regimes ranging from weak-field orbits to strong-field photon capture.

\begin{itemize}

% ============================================================
% 1. NULL GEODESIC NORMALIZATION
% IMAGE LEFT -- TEXT RIGHT
% ============================================================
\newpage
\item \textbf{Null Geodesic Normalization.}

\vspace{2mm}

\noindent
\parbox[t]{0.32\linewidth}{
\vspace{0pt}
\centering
\refstepcounter{figure}
\includegraphics[width=\linewidth]{null geodesic.jpeg}

\vspace{1mm}

\footnotesize
\textbf{Figure~\thefigure:}
Null geodesic normalization residual
$g_{\mu\nu}k^{\mu}k^{\nu}$ tracked over the course of integration,
remaining within $\pm10^{-14}$ for the majority of the trajectory.
\label{fig:null-normalization}
}
\hfill
\parbox[t]{0.64\linewidth}{
\vspace{0pt}
In addition to the timelike conservation checks described previously,
the null condition
\begin{equation}
g_{\mu\nu}k^{\mu}k^{\nu}=0
\end{equation}
governing photon trajectories was monitored at every integration step.
As shown in Figure~\ref{fig:null-normalization}, the normalization
residual remains bounded to order $10^{-14}$ across the majority of
the trajectory, consistent with double-precision floating-point
accuracy, with a modest growth in residual magnitude observed only
in the high-curvature regime near closest approach --- an expected
consequence of increased sensitivity in the Christoffel symbols near
the photon sphere rather than a defect in the integration scheme.
}

\vspace{5mm}


% ============================================================
% 2. EFFECTIVE POTENTIAL
% TEXT LEFT -- IMAGE RIGHT
% ============================================================

\item \textbf{Effective Potential and Characteristic Radii.}

\vspace{2mm}

\noindent
\parbox[t]{0.64\linewidth}{
\vspace{0pt}
The Schwarzschild effective potential for massive test particles was
computed and compared against its known analytic features: the event
horizon at $r=2M$, the photon sphere at $r=3M$, and the innermost
stable circular orbit (ISCO) at $r=6M$. As shown in
Figure~\ref{fig:effective-potential}, these radii are included as reference markers. For massive particles, the limit of the effective potential depend on the particle's angular momentum and therefore does not coincide with the photon-sphere radius and produces the
inflection behavior at the ISCO, confirming that the numerical
framework recovers the correct qualitative and quantitative structure
of particle dynamics predicted analytically for Schwarzschild
spacetime.
}
\hfill
\parbox[t]{0.32\linewidth}{
\vspace{0pt}
\centering
\refstepcounter{figure}
\includegraphics[width=\linewidth]{effective potential.jpeg}

\vspace{1mm}

\footnotesize
\textbf{Figure~\thefigure:}
Schwarzschild effective potential for a massive test particle,
with the event horizon, photon sphere, and ISCO marked at their
analytically predicted radii.
\label{fig:effective-potential}
}

\vspace{5mm}


% ============================================================
% 3. TRAJECTORY CLASSES
% IMAGE LEFT -- TEXT RIGHT
% ============================================================
\newpage
\item \textbf{Classification of Particle Trajectory Regimes.}

\vspace{2mm}

\noindent
\parbox[t]{0.32\linewidth}{
\vspace{0pt}
\centering
\refstepcounter{figure}
\includegraphics[width=\linewidth]{multiple particle simulation.jpeg}

\vspace{1mm}

\footnotesize
\textbf{Figure~\thefigure:}
Superposed particle trajectories for the four characteristic
regimes predicted by the effective potential --- circular,
elliptical, escape, and radial infall --- around a Schwarzschild
black hole.
\label{fig:trajectory-classes}
}
\hfill
\parbox[t]{0.64\linewidth}{
\vspace{0pt}
To validate that the geodesic solver correctly reproduces the four
qualitatively distinct trajectory classes predicted by the effective
potential, test particles were initialized with energies and angular
momenta corresponding to a stable circular orbit, a bound elliptical
orbit, a marginally unbound escape trajectory, and a purely radial
infall. The resulting trajectories, shown together in
Figure~\ref{fig:trajectory-classes}, correctly reproduce the
expected qualitative behavior for each regime, with the escape and
infall cases terminating at the event horizon or diverging to large
$r$ as expected.
}

\vspace{5mm}


% ============================================================
% 4. PERIHELION PRECESSION
% TEXT LEFT -- IMAGES RIGHT
% ============================================================

\item \textbf{Relativistic Perihelion Precession.}

\vspace{2mm}

\noindent
\begin{minipage}{\linewidth}
\centering

\refstepcounter{figure}
\begin{minipage}[t]{0.46\linewidth}
\centering
\includegraphics[width=0.95\linewidth]{epiltical orbit.jpeg}

\vspace{1mm}

\footnotesize
\textbf{Figure~\thefigure:}
Bound orbit integrated over multiple radial periods, showing
progressive advance of the periapsis.
\label{fig:precession-advance}
\end{minipage}
\hfill
\refstepcounter{figure}
\begin{minipage}[t]{0.46\linewidth}
\centering
\includegraphics[width=\linewidth]{relatavistic orbit.jpeg}

\vspace{1mm}

\footnotesize
\textbf{Figure~\thefigure:}
Long-duration integration of a bound relativistic orbit,
tracing a non-closing rosette pattern characteristic of
strong-field perihelion precession.
\label{fig:precession-rosette}
\end{minipage}

\vspace{4mm}

\noindent
Unlike the closed, non-precessing ellipses predicted by Newtonian
gravity, bound orbits in the Schwarzschild metric exhibit
perihelion precession --- a hallmark signature of general
relativistic dynamics. This effect was validated by integrating
a bound orbit over multiple radial periods and confirming that the
orbit fails to close, instead advancing its periapsis with each
revolution (Figure~\ref{fig:precession-advance}) and tracing a
rosette pattern over many periods
(Figure~\ref{fig:precession-rosette}), qualitatively consistent
with the precession observed in relativistic orbits such as
Mercury's in the weak-field limit and S2's orbit around Sgr A*
in the strong-field regime.

\end{minipage}

\vspace{5mm}

\newpage
% ============================================================
% 5. WEAK FIELD LIMIT
% IMAGE LEFT -- TEXT RIGHT
% ============================================================

\item \textbf{Weak-Field (Newtonian) Limit Recovery.}

\vspace{2mm}

\noindent
\parbox[c]{0.32\linewidth}{
\centering
\refstepcounter{figure}
\includegraphics[width=\linewidth]{partical orbit.jpeg}

\vspace{1mm}

\footnotesize
\textbf{Figure~\thefigure:}
Circular orbit at large radial distance from the black hole,
correctly closing on itself and recovering the Newtonian
weak-field limit.
\label{fig:weak-field}
}
\hfill
\parbox[c]{0.64\linewidth}{
As a consistency check on the opposite limit, a test particle was
placed on a stable circular orbit at a radial distance far exceeding
the ISCO, where relativistic corrections are expected to vanish.
The resulting trajectory (Figure~\ref{fig:weak-field}) traces a
closed, non-precessing circular orbit, confirming that the geodesic
integrator correctly recovers Newtonian circular-orbit behavior
in the weak-field limit.
}

\vspace{5mm}


% ============================================================
% 6. PHOTON DEFLECTION
% TEXT LEFT -- IMAGE RIGHT
% ============================================================

\item \textbf{Photon Deflection and the Critical Impact Parameter.}

\vspace{2mm}

\noindent
\parbox[t]{0.64\linewidth}{
\vspace{0pt}
The gravitational deflection of light was validated by tracing
photon trajectories at several impact parameters $b$ relative to
the analytically derived critical impact parameter

\begin{equation}
b_{\text{crit}}=3\sqrt{3}\,M\approx5.196\,M.
\end{equation}

As shown in Figure~\ref{fig:photon-deflection}, rays with impact
parameters below $b_{\text{crit}}$ are captured by the black hole, while those at $b=5$, $6$, and $8$ escape with progressively
smaller net deflection angles as $b$ increases consistent with
the analytic prediction validated separately in
\texttt{validate\_critical\_impact\_parameter.py}.
}
\hfill
\parbox[t]{0.32\linewidth}{
\vspace{0pt}
\centering
\refstepcounter{figure}
\includegraphics[width=\linewidth]{light bending.jpeg}

\vspace{1mm}

\footnotesize
\textbf{Figure~\thefigure:}
Photon trajectories at increasing impact parameter $b$, showing
progressively reduced gravitational deflection as $b$ increases
beyond $b_{\text{crit}}$.
\label{fig:photon-deflection}
}

\vspace{5mm}

\newpage
% ============================================================
% 7. KERR ERGOSPHERE
% IMAGE LEFT -- TEXT RIGHT
% ============================================================

\item \textbf{Kerr Ergosphere Geometry.}

\vspace{2mm}

\noindent
\parbox[t]{0.32\linewidth}{
\vspace{0pt}
\centering
\refstepcounter{figure}
\includegraphics[width=\linewidth]{ergo sphere.jpeg}

\vspace{1mm}

\footnotesize
\textbf{Figure~\thefigure:}
Ergosphere and event horizon geometry for a Kerr black hole with
spin parameter $a=0.9M$, showing the expected oblate ergosphere
touching the horizon at the poles.
\label{fig:kerr-ergosphere}
}
\hfill
\parbox[t]{0.64\linewidth}{
\vspace{0pt}
For the rotating (Kerr) case, the spatial extent of the ergosphere
was validated against its known analytic form,

\begin{equation}
r_E(\theta)
=
M+\sqrt{M^2-a^2\cos^2\theta}.
\end{equation}

which is oblate and touches the event horizon only at the poles
($\theta=0,\pi$) while extending beyond it at the equator.
Figure~\ref{fig:kerr-ergosphere} ($a=0.9M$) confirms this expected
geometry, with the ergosphere visibly enclosing the event horizon
and coinciding with it only along the rotation axis, validating
the frame-dragging structure introduced by nonzero spin.
}

\vspace{5mm}


\end{itemize}


Collectively, these validations span both the null and timelike
geodesic sectors and cover the full range of dynamical regimes relevant
to the rendering pipeline --- from photon capture and escape, through
bound and marginally bound massive-particle orbits, to the rotating
Kerr case --- providing strong confidence in the physical correctness
of the underlying simulation before it is used to generate training
data for the surrogate model described in the following section.

% ============================================================
% MACHINE LEARNING SURROGATE MODEL
% ============================================================

\section{Machine Learning Surrogate Model for Photon Geodesics}

To reduce the computational cost associated with numerical integration of photon geodesics, a supervised machine learning framework is developed as a surrogate model for Schwarzschild ray tracing. The objective of this model is to learn the nonlinear mapping between the initial ray parameters and the corresponding physical observables obtained from the validated, high-precision numerical geodesic solver described in the preceding sections. A dataset is first generated by solving the geodesic equations over a wide range of impact parameters, thereby providing ground-truth solutions for training and evaluation.

The input feature vector consists of the photon impact parameter ($b$) and the normalized impact parameter ($b/b_{\text{crit}}$), where $b_{\text{crit}}$ denotes the critical impact parameter corresponding to the photon sphere. The prediction targets include the minimum radial distance ($r_{\min}$), the final azimuthal angle ($\varphi_{\text{final}}$), the gravitational deflection angle ($\alpha$), and the photon capture state. Since these outputs represent both continuous and categorical quantities, the problem is formulated as a combination of regression and binary classification tasks.

\begin{figure}[htbp]
    \centering

    % Top row
    \begin{subfigure}{0.48\linewidth}
        \centering
        \includegraphics[width=\linewidth]{WhatsApp Image 2026-08-10 at 5.25.47 PM.jpeg}
        \caption{$r_{\min}$}
    \end{subfigure}
    \hfill
    \begin{subfigure}{0.48\linewidth}
        \centering
        \includegraphics[width=\linewidth]{WhatsApp Image 2026-08-08 at 11.59.28 PM.jpeg}
        \caption{Predicted}
    \end{subfigure}

    \vspace{0.5cm}

    % Bottom row
    \begin{subfigure}{0.48\linewidth}
        \centering
        \includegraphics[width=\linewidth]{WhatsApp Image 2026-08-08 at 11.58.47 PM.jpeg}
        \caption{Binary}
    \end{subfigure}
    \hfill
    \begin{subfigure}{0.48\linewidth}
        \centering
        \includegraphics[width=\linewidth]{WhatsApp Image 2026-08-08 at 11.57.56 PM.jpeg}
        \caption{Fourth Image}
    \end{subfigure}

    \caption{Comparison of the predicted and corresponding ray-tracing outputs.}
    \label{fig:comparison}
\end{figure}

The dataset is randomly partitioned into training (75\%), validation (10\%), and testing (15\%) subsets. Prior to training, the input features are standardized using a single feature scaler fitted on the training data, while independent target scalers are employed for each regression output to account for differences in their numerical distributions. Initially, separate single-task neural networks are trained for each target variable to establish baseline performance and evaluate the predictive capability of the proposed surrogate for individual physical quantities. Subsequently, a multi-task learning architecture is constructed by introducing shared hidden layers for common feature extraction, together with task-specific output heads for each prediction target. This architecture enables the model to exploit shared physical relationships among the different observables while retaining dedicated prediction branches for the regression and classification tasks.

The neural networks are implemented as fully connected feed-forward architectures and optimized using the Adam optimizer, with Mean Squared Error (MSE) loss for the regression tasks and Binary Cross-Entropy (BCE) loss for the photon-capture classification task. Early stopping based on validation loss is employed throughout training to mitigate overfitting. Model performance is evaluated using Mean Absolute Error (MAE), Root Mean Squared Error (RMSE), and the coefficient of determination ($R^2$) for the regression outputs, while classification performance is assessed using accuracy, precision, recall, and F1-score. The predictive accuracy of the proposed surrogate model is finally compared against the numerical geodesic solver to evaluate its suitability as an approximation framework for real-time computational ray tracing.

\begin{table}[htbp]
\centering
\caption{Performance of the Multitask Regression Model}
\label{tab:multitask_regression}
\begin{tabular}{lc}
\hline
\textbf{Target Variable} & \textbf{$R^2$ Score} \\
\hline
$r_{\min}$ & 0.9999010170 \\
$\phi$    & 0.9928728075 \\
$\alpha$  & 0.9955208800 \\
\hline
\end{tabular}
\end{table}

\begin{table}[htbp]
\centering
\caption{Binary Classification Performance}
\label{tab:classification}
\begin{tabular}{lcccc}
\hline
\textbf{Class} & \textbf{Precision} & \textbf{Recall} & \textbf{F1-Score} & \textbf{Support} \\
\hline
0 & 1.00 & 0.93 & 0.96 & 1043 \\
1 & 0.48 & 1.00 & 0.65 & 69 \\
\hline
\textbf{Accuracy} & & & \textbf{0.93} & 1112 \\
\textbf{Macro Avg.} & 0.74 & 0.96 & 0.81 & 1112 \\
\textbf{Weighted Avg.} & 0.97 & 0.93 & 0.94 & 1112 \\
\hline
\end{tabular}
\end{table}

% ============================================================
% MULTI-TASK ARCHITECTURE
% ============================================================

\subsection{Multi-Task Neural Network Architecture}

The overall workflow of the proposed multi-task neural network is illustrated in Figure~\ref{fig:mtl-architecture}. The model consists of a shared feature-extraction backbone --- three fully connected layers of widths 64, 128, and 64, each followed by a ReLU activation --- followed by four task-specific prediction heads corresponding to the four physical observables of interest. Each head comprises two additional dense layers (widths 32 and 16) that specialize the shared representation toward its respective target, before producing a final scalar output: three continuous regression outputs ($r_{\min}$, $\varphi_{\text{final}}$, and $\alpha$) and one binary classification output (captured).

\begin{figure}[htbp]
    \centering

    \begin{minipage}[t]{0.48\textwidth}
        \centering
        \includegraphics[width=\linewidth]{Figure_1.png}
        \caption{PINN vs Data driven model}
        \label{fig:pinn_vs_data}
    \end{minipage}
    \hfill
    \begin{minipage}[t]{0.48\textwidth}
        \centering
        \includegraphics[width=\linewidth]{Figure_2.png}
        \caption{Data loss vs Physics Constraint Convergence}
        \label{fig:data_vs_physics}
    \end{minipage}

\end{figure}

This shared-backbone, multi-head design reflects the underlying physical structure of the problem: all four observables are deterministic functions of the same impact parameter and are governed by the same underlying geodesic equation, motivating a shared representation that can, in principle, be learned jointly rather than separately across four independent networks. The task-specific heads then allow each output branch to adapt this shared representation to its own numerical scale and loss function — regression heads trained under MSE loss and the classification head trained under BCE loss, without requiring the shared layers to compromise between conflicting objectives.

\begin{figure}[htbp]
    \centering
    \includegraphics[width=0.45\textwidth]{image.png}
    \caption{Architecture of the proposed multi-task neural network surrogate model. Shared hidden layers learn common physical representations, while task-specific branches predict the individual photon trajectory observables.}
    \label{fig:mtl-architecture}
\end{figure}
 \newpage
\section{Appendix}
\subsection{\textbf{Some definitions:}}
\subsubsection{\textbf{Event Horizon:}} The event horizon is the boundary surrounding a black hole beyond which no matter, light, or information can escape to distant observers. It represents the limit of causal communication with the external universe.
\subsubsection{\textbf{Innermost Stable Circular Orbit(ISCO):}} The ISCO is the smallest radius at which a massive particle can maintain a stable circular orbit around a black hole. Inside the ISCO, any small inward perturbation causes the particle to leave its circular orbit and plunge toward the black hole's event horizon.
\subsubsection{\textbf{Photon Sphere:}} The photon sphere is the radius at which null geodesics admit circular solutions. These circular photon orbits are unstable, so even a small perturbation causes the photon either to escape to infinity or to spiral into the black hole. 
\subsubsection{\textbf{Impact parameter(b):}} The distance between an incoming light ray's straight-line path and a parallel line passing through the centre of the black hole. It is equal to the ratio of angular momentum to the energy of a photon.
\begin{equation}
b= \frac{L}{E}
\end{equation}

\subsubsection{\textbf{Einstein Ring:}} An Einstein Ring, also known as an Einstein-Chwolson ring, is created when light from a galaxy or star passes by a massive object en route to Earth. Due to gravitational lensing, the light is diverted, making it seem to come from different places.  
\subsubsection{\textbf{Singularity:}} The singularity is the point where the mathematical description breaks down (diverges).
We get two types of singularities in a spherical coordinate system:
\\
\underline{Coordinate Singularity:} The metric coefficient becomes 0 or infinite, but the curvature remains finite, and a different coordinate system removes the apparent singularity. 
\\
\underline{Curvature Singularity(r=0):} Curvature invariants, such as Kretschmann scalar, diverge. No coordinate transformation removes this.

\subsubsection{Affine Parameter($\lambda$):} An affine parameter is a parameter used to label the points along a geodesic in a way that preserves the natural uniform progression along the path. For massive particles, proper time can serve as such a parameter. However, photons follow null geodesics and experience no proper time, so proper time cannot be used to parameterize their trajectories. Instead, an affine parameter $\lambda$ is introduced as a mathematical measure of progress along the null geodesic.




\newpage\section{Acknowledgement}
We would like to express our heartfelt gratitude to the IAPT (Indian Association of Physics Teachers) for providing us with the opportunity and platform to present this project. This project could not have been possible without the help of the University of Allahabad, which served as our primary workspace throughout the course of this project. 

We would also like to thank our professors, who offered valuable insights to help with the completion of this project. 
We are incredibly grateful to our friends and cohort members for their editing help, feedback sessions, and moral support. 

Furthermore, we owe the deepest gratitude to all the pioneers of this field, for we could never have brought this project into existence without all the people who have dedicated their lives to the concepts this project covers. As {Sir Isaac Newton} very rightfully said, "If I have seen further, it is by standing on the shoulders of giants".
\newpage
\begin{thebibliography}{99}   
\bibitem{mitfieldeqs}
M. Tegmark,
\textit{Einstein's Field Equations},
Physics 8.033: Relativity,
Massachusetts Institute of Technology,
2002.
\cite{mitfieldeqs}
\bibitem{einstein1920}
A. Einstein,
\textit{Relativity: The Special and General Theory},
translated by R. W. Lawson,
Methuen \& Co., London, 1920.
\cite{einstein1920}
\bibitem{hawking2002}
Hawking, S. W., \textit{On the Shoulders of Giants:
The Great Works of Physics and Astronomy}, Running Press.
\cite{hawking2002}
\bibitem{callahan2000}
J. J. Callahan,
\textit{The Geometry of Spacetime: An Introduction to Special and General Relativity},
Springer, New York, 2000.
\cite{callahan2000}
\bibitem{luminet1979}
J.-P. Luminet,
\textit{Image of a spherical black hole with thin accretion disk},
Astronomy \& Astrophysics, vol. 75, pp. 228--235, 1979.
\cite{luminet1979}
\bibitem{schwarzschild1916}
K. Schwarzschild,
\textit{Über das Gravitationsfeld eines Massenpunktes nach der Einsteinschen Theorie},
Sitzungsberichte der Königlich Preussischen Akademie der Wissenschaften zu Berlin,
pp. 189--196, 1916.
\cite{schwarzschild1916}
\bibitem{kerr1963}
R. P. Kerr,
\textit{Gravitational Field of a Spinning Mass as an Example of Algebraically Special Metrics},
Physical Review Letters, vol. 11, pp. 237--238, 1963.
\cite{kerr1963}
\bibitem{sun2025}
M. Sun, Z. Fang, J. Wang, K. Zhang, Q. Zhang, and R. Xu,
\textit{Learning Null Geodesics for Gravitational Lensing Rendering in General Relativity},
arXiv:2507.15775 [gr-qc], 2025.
\cite{sun2025}
\bibitem{james2015}
O. James, E. von Tunzelmann, P. Franklin, and K. S. Thorne,
\textit{Gravitational Lensing by Spinning Black Holes in Astrophysics,
and in the Movie Interstellar},
Classical and Quantum Gravity, vol. 32, no. 6, 065001, 2015.
\cite{james2015}
\bibitem{dexter2009}
J. Dexter and E. Agol,
\textit{A Fast New Public Code for Computing Photon Orbits in a Kerr Spacetime},
The Astrophysical Journal, vol. 696, no. 2, pp. 1616--1629, 2009.
\cite{dexter2009}
\bibitem{vincent2011}
F. H. Vincent, T. Paumard, E. Gourgoulhon, and G. Perrin,
\textit{GYOTO: a new general relativistic ray-tracing code},
Classical and Quantum Gravity, vol. 28, no. 22, 225011, 2011.
\cite{vincent2011}
\bibitem{chan2013}
C.-K. Chan, D. Psaltis, F. Özel, R. Narayan, A. S. Marrone,
and others,
\textit{GRay: a Massively Parallel GPU-Based Code for Ray Tracing in Relativistic Spacetimes},
The Astrophysical Journal, vol. 777, no. 1, 13, 2013.
\cite{chan2013}
\bibitem{eyama2025}
S. Eyama and Y. Masada,
\textit{Constraining the Cosmological Constant from Stellar Orbits Around Sgr A* Using Physics-Informed Neural Networks},
arXiv:2508.19719 [astro-ph.CO], 2025.
\cite{eyama2025}
\bibitem{raissi2019}
M. Raissi, P. Perdikaris, and G. E. Karniadakis,
\textit{Physics-informed neural networks: A deep learning framework for
solving forward and inverse problems involving nonlinear partial
differential equations},
Journal of Computational Physics, vol. 378, pp. 686--707, 2019.
doi:10.1016/j.jcp.2018.10.045.
\cite{raissi2019}
\end{thebibliography}
\end{itemize}
\end{document}
